% ===== §31 =====
\section{The Ethics of Cultural Cognition and Practice}
\label{sec:ethics}

\noindent
The turn of the previous section, from the diagnosis of capture to the conditions of recovery, opens onto a normative question the paper must now face directly. If the capacity for endogenous generation is never wholly destroyed, and if the reclamation of a captured culture is possible, then the question arises what it would mean to know and to practise intimate culture rightly---what ethics attaches to the cognition and the practice of a culture once it is understood as the endogenous, generative sediment the paper has shown it to be. This section states that ethics, in two layers, the cognitive and the practical, and shows that they are, at bottom, one.

\subsection*{The cognitive layer: culture is not an object to be possessed}

The ethics begins in cognition, because how one knows intimate culture already carries a normative weight. The paper has shown that intimate culture is not a content its forms contain but an activity its participants perform, known only from within, in the participatory circuit that its meaning consists in. From this it follows that to \emph{know} an intimate culture rightly is not to objectify it---not to hold it at the distance of a thing surveyed, inventoried, possessed---but to participate in the generative activity that is its life. And this cognitive fact is already an ethical one, for the misrecognition that treats a living, generated culture as a possessable object is not merely an error but a failure of the right relation to the culture and to the other with whom it is shared. To treat the bond's culture as a thing one has is to have already withdrawn from the participation in which alone it lives; to treat it as an asset one can exhibit and convert is to have already delivered it toward the capture that draws its substance off. The cognitive error and the ethical failure are the same act seen twice: mistaking the generated for the given, the activity for the object, the participation for the possession. The first requirement of the ethics is therefore cognitive: to know intimate culture as what it is, a generation to be participated in, and not as what the misrecognitions would make it, a thing to be had or shown.

\subsection*{The practical layer: non-possessive generation}

From the cognitive requirement follows the practical one, and it can be stated as a single principle: the practice proper to intimate culture is a \emph{non-possessive generation}. The culture is to be generated and re-generated, continually, in the participatory activity that is its life; and it is to be generated without being possessed---without being seized by either party as a thing they own, without being extracted as an asset, without being turned into an instrument of one party's power over the other. Non-possessive generation is the practice in which the culture is made and kept alive precisely by not being grasped as a made thing: made continually, held by neither, owned by no one, because its mode of being is the activity of its generation and not the possession of its product. The classical tradition of this series offers, as one resource among others for naming this practice, the conception of a generation that does not possess what it generates---that produces without owning, acts without holding onto, brings to maturity without ruling over---and the reader familiar with that tradition will recognize the resonance; but the principle stands on the argument of this paper, and does not depend on the classical formulation, which is invoked here as an illumination and not as a ground. What the argument requires is only this: that intimate culture, being a generation and not a possession, is rightly practised only in a generation that refuses to possess, and that every grasping of it as a possession---whether toward private hoarding or toward public exhibition---is a lapse from the practice its being demands.

\subsection*{The power dimension: non-possession and non-entrenchment are one posture}

This practical ethics joins, at its root, the ethics of power the synthesis established, and marking the junction shows that the paper's normative claims form a single posture rather than a list. The sublation of power held that a just intimate culture must generate, alongside the power its generation carries, the counter-power that prevents that power from entrenching---the friction against accumulation and the force of re-orientation. The ethics of non-possessive generation is the same requirement seen from another side. To possess an intimate culture---to grasp it as a thing one owns---is precisely to entrench a power in it, to fix the generative process into a possessed structure that one party holds and the other is held by; and to refuse possession, to keep the culture in the condition of ongoing generation owned by neither, is precisely to prevent that entrenchment, to keep the power in the culture in the condition of power-to rather than power-over. Non-possession and non-entrenchment are therefore not two requirements but one posture, seen in the register of ownership and in the register of power respectively: the refusal to possess the culture is the refusal to entrench the power, and the practice of non-possessive generation is the practice of a culture in which power remains enabling rather than dominating. The two ethical claims the paper has made, the one about ownership and the one about power, are one claim about the single posture proper to a generative intimacy.

\subsection*{Against cultural capital, once more}

The ethics can be sharpened by setting it against the concept it has opposed since Part One. Bourdieu's cultural capital is the very type of a possessed and entrenched culture: it is accrued, held, converted, deployed as an asset in a field of struggle, and its whole being is possession and the power that possession confers \citep{bourdieu1986}. The ethics of non-possessive generation is the exact negation of this: a culture whose worth is not in its possession but in its generation, not in the advantage it confers on its holder but in the world it opens between its participants, not extractable because not possessable, not a capital because not a thing. The contrast is not merely conceptual but ethical, for it marks two ways of standing toward what one shares with another: the way of the asset, which holds, converts, and extracts, and the way of the generation, which makes, participates, and lets be. The whole of this paper's ethics can be compressed into the choice between these two, and the choice is not neutral. To practise intimate culture as cultural capital is to have already begun to capture it oneself, to stand toward one's own bond as the market stands toward it, holding its meaning as an asset to be converted; to practise it as non-possessive generation is to keep it in the register of the relational Real that no market can reach, generated and re-generated between two who own it by not owning it. The reclamation the previous section showed to be possible is, in practice, exactly this: the return, from the possession and exhibition that deliver intimate culture to capture, to the non-possessive generation in which it lives.
